
\section{Cursos}

En esta sección se incluye la lista de cursos de la propuesta, organizada por año y, cuando corresponde, por énfasis. Cada entrada enlaza al programa recogido en la sección~\ref{sec:anexo-programas}.

\subsection{Cursos de primer año}

La siguiente tabla muestra los niveles de dominio y contenidos esperados para estos dos semestres.

\begin{table}[H]
\begin{center}
\begin{tabular}{|ccc|}
\hline
\multicolumn{3}{|c|}{\cellcolor[HTML]{00C0F3} \textbf{Niveles de dominio y contenidos por área curricular}} \\ \hline
\multicolumn{1}{|l|}{\cellcolor[HTML]{8ED8F8}\textbf{Área curricular}} & \multicolumn{1}{l|}{\cellcolor[HTML]{8ED8F8}\textbf{Dominio}} & \cellcolor[HTML]{8ED8F8}\textbf{Contenidos} \\ \hline
\multicolumn{1}{|c|}{Habilidades operacionales} & \multicolumn{1}{c|}{Alto} & Alto \\ \hline
\multicolumn{1}{|l|}{Pensamiento abstracto} & \multicolumn{1}{c|}{Medio} & Bajo \\ \hline
\multicolumn{1}{|l|}{Pensamiento algorítmico} & \multicolumn{1}{c|}{Bajo} & Bajo \\ \hline
\multicolumn{1}{|l|}{Formalismo} & \multicolumn{1}{c|}{Bajo} & Bajo \\ \hline
\multicolumn{1}{|c|}{Habilidades investigativas} & \multicolumn{1}{c|}{-} & Bajo \\ \hline
\end{tabular}
\end{center}
\caption{Niveles de dominio y contenidos para los cursos del primer año de la carrera.}
\end{table}

\subsubsection*{Cursos comunes a los dos énfasis}

\begin{itemize}
    \item \hyperlink{MA0151}{MA-0151 Fundamentos de álgebra, trigonometría y geometría analítica}
    \item \hyperlink{MA0152}{MA-0152 Matemática Exploratoria}
    \item \hyperlink{MA0251}{MA-0251 Cálculo en una variable}
    \item \hyperlink{MA0252}{MA-0252 Introducción a las Demostraciones}
\end{itemize}


\subsection{Cursos de segundo año}

\begin{table}[H]
\begin{center}
\begin{tabular}{|ccc|}
\hline
\multicolumn{3}{|c|}{\cellcolor[HTML]{00C0F3} \textbf{Niveles de dominio y contenidos por área curricular}} \\ \hline
\multicolumn{1}{|l|}{\cellcolor[HTML]{8ED8F8}\textbf{Área curricular}} & \multicolumn{1}{l|}{\cellcolor[HTML]{8ED8F8}\textbf{Dominio}} & \cellcolor[HTML]{8ED8F8}\textbf{Contenidos} \\ \hline
\multicolumn{1}{|l|}{Habilidades operacionales} & \multicolumn{1}{c|}{Alto} & Medio \\ \hline
\multicolumn{1}{|l|}{Pensamiento abstracto} & \multicolumn{1}{c|}{Medio} & Medio \\ \hline
\multicolumn{1}{|l|}{Pensamiento algorítmico} & \multicolumn{1}{c|}{Medio} & Medio \\ \hline
\multicolumn{1}{|l|}{Formalismo} & \multicolumn{1}{c|}{Medio} & Medio \\ \hline
\multicolumn{1}{|l|}{Habilidades investigativas} & \multicolumn{1}{c|}{-} & Bajo \\ \hline
\end{tabular}
\end{center}
\caption{Niveles de dominio y contenidos para los cursos del segundo año de la carrera.}
\end{table}

\subsubsection*{Cursos comunes a los dos énfasis}

\begin{itemize}
    \item \hyperlink{MA0341}{MA-0341 Matemática Computacional}
    \item \hyperlink{MA0351}{MA-0351 Cálculo en varias variables}
    \item \hyperlink{MA0361}{MA-0361 Algebra Lineal I}
    \item \hyperlink{MA0451}{MA-0451 Principios de análisis I}
    \item \hyperlink{MA0461}{MA-0461 Algebra Lineal II}
\end{itemize}

\subsubsection*{Cursos del énfasis en matemática pura}
\begin{itemize}
    \item \hyperlink{MA0471}{MA-0471 Introducción a la Geometría Diferencial}
    \item \hyperlink{MA0496}{MA-0496 Teoría de Números}
\end{itemize}

\subsubsection*{Cursos del énfasis en matemática aplicada}
\begin{itemize}
    \item {CA-0204 Herramientas de ciencia de datos I}
    \item \hyperlink{CA0721}{CA-0721 Probabilidad}
\end{itemize}




\subsection{Cursos de tercer año}

La siguiente tabla resume los niveles de dominio y contenidos esperados para el tercer año, según las áreas curriculares definidas en este capítulo.

\begin{table}[H]
\begin{center}
\begin{tabular}{|ccc|}
\hline
\multicolumn{3}{|c|}{\cellcolor[HTML]{00C0F3} \textbf{Niveles de dominio y contenidos por área curricular}} \\ \hline
\multicolumn{1}{|l|}{\cellcolor[HTML]{8ED8F8}\textbf{Área curricular}} & \multicolumn{1}{l|}{\cellcolor[HTML]{8ED8F8}\textbf{Dominio}} & \cellcolor[HTML]{8ED8F8}\textbf{Contenidos} \\ \hline
\multicolumn{1}{|l|}{Habilidades operacionales} & \multicolumn{1}{c|}{Alto} & Bajo \\ \hline
\multicolumn{1}{|l|}{Pensamiento abstracto} & \multicolumn{1}{c|}{Alto} & Alto \\ \hline
\multicolumn{1}{|l|}{Pensamiento algorítmico} & \multicolumn{1}{c|}{Alto} & Alto \\ \hline
\multicolumn{1}{|l|}{Formalismo} & \multicolumn{1}{c|}{Alto} & Medio \\ \hline
\multicolumn{1}{|l|}{Habilidades investigativas} & \multicolumn{1}{c|}{Bajo} & Bajo \\ \hline
\end{tabular}
\end{center}
\caption{Niveles de dominio y contenidos para los cursos del tercer año de la carrera.}
\end{table}

\subsubsection*{Cursos comunes a los dos énfasis}

\begin{itemize}
    \item \hyperlink{MA0515}{MA-0515 Ecuaciones Diferenciales Ordinarias}
    \item \hyperlink{MA0551}{MA-0551 Principios de Análisis II}
    \item \hyperlink{MA0635}{MA-0635 Topología General}
    \item \hyperlink{MA0641}{MA-0641 Análisis Numérico I}
\end{itemize}

\subsubsection*{Cursos del énfasis en matemática pura}
\begin{itemize}
    \item \hyperlink{MA0561}{MA-0561 Algebra Abstracta I}
    \item \hyperlink{MA0661}{MA-0661 Algebra Abstracta II}
\end{itemize}

\subsubsection*{Cursos del énfasis en matemática aplicada}
\begin{itemize}
    \item {CA-0303 Estadística Actuarial I}
    \item \hyperlink{MA0647}{MA-0647 Modelación Matemática}
\end{itemize}

\subsection{Cursos de cuarto año}

La siguiente tabla resume los niveles de dominio y contenidos esperados para el cuarto año, según las áreas curriculares definidas en este capítulo.

\begin{table}[H]
\begin{center}
\begin{tabular}{|ccc|}
\hline
\multicolumn{3}{|c|}{\cellcolor[HTML]{00C0F3} \textbf{Niveles de dominio y contenidos por área curricular}} \\ \hline
\multicolumn{1}{|l|}{\cellcolor[HTML]{8ED8F8}\textbf{Área curricular}} & \multicolumn{1}{l|}{\cellcolor[HTML]{8ED8F8}\textbf{Dominio}} & \cellcolor[HTML]{8ED8F8}\textbf{Contenidos} \\ \hline
\multicolumn{1}{|l|}{Habilidades operacionales} & \multicolumn{1}{c|}{Alto} & Bajo \\ \hline
\multicolumn{1}{|l|}{Pensamiento abstracto} & \multicolumn{1}{c|}{Alto} & Alto \\ \hline
\multicolumn{1}{|l|}{Pensamiento algorítmico} & \multicolumn{1}{c|}{Alto} & Alto \\ \hline
\multicolumn{1}{|l|}{Formalismo} & \multicolumn{1}{c|}{Alto} & Alto \\ \hline
\multicolumn{1}{|l|}{Habilidades investigativas} & \multicolumn{1}{c|}{Medio} & Medio \\ \hline
\end{tabular}
\end{center}
\caption{Niveles de dominio y contenidos para los cursos del cuarto año de la carrera.}
\end{table}

\subsubsection*{Cursos comunes a los dos énfasis}

\begin{itemize}
    \item \hyperlink{MA0702}{MA-0702 Análisis Complejo}
    \item \hyperlink{MA0705}{MA-0705 Análisis Real I}
    \item \hyperlink{MA0880}{MA-0880 Comunicación en las ciencias}
\end{itemize}

\subsubsection*{Cursos del énfasis en matemática pura}
\begin{itemize}
    \item \hyperlink{MA0881Pura}{MA-0881 Práctica profesional en matemática pura}
\end{itemize}

\subsubsection*{Cursos del énfasis en matemática aplicada}
\begin{itemize}
    \item {CA-0411 Análisis de Datos I}
    \item \hyperlink{MA0883}{MA-0883 Práctica profesional en matemática aplicada}
\end{itemize}

\subsection{Cursos de licenciatura}

Las áreas curriculares del capítulo de estructura curricular contemplan explícitamente los años I a IV; para la licenciatura se proyecta la culminación de esa progresión, con énfasis reforzado en habilidades investigativas mediante los seminarios de estudios dirigidos.

\begin{table}[H]
\begin{center}
\begin{tabular}{|ccc|}
\hline
\multicolumn{3}{|c|}{\cellcolor[HTML]{00C0F3} \textbf{Niveles de dominio y contenidos por área curricular}} \\ \hline
\multicolumn{1}{|l|}{\cellcolor[HTML]{8ED8F8}\textbf{Área curricular}} & \multicolumn{1}{l|}{\cellcolor[HTML]{8ED8F8}\textbf{Dominio}} & \cellcolor[HTML]{8ED8F8}\textbf{Contenidos} \\ \hline
\multicolumn{1}{|l|}{Habilidades operacionales} & \multicolumn{1}{c|}{Alto} & Bajo \\ \hline
\multicolumn{1}{|l|}{Pensamiento abstracto} & \multicolumn{1}{c|}{Alto} & Alto \\ \hline
\multicolumn{1}{|l|}{Pensamiento algorítmico} & \multicolumn{1}{c|}{Alto} & Alto \\ \hline
\multicolumn{1}{|l|}{Formalismo} & \multicolumn{1}{c|}{Alto} & Alto \\ \hline
\multicolumn{1}{|l|}{Habilidades investigativas} & \multicolumn{1}{c|}{Alto} & Alto \\ \hline
\end{tabular}
\end{center}
\caption{Niveles de dominio y contenidos proyectados para los cursos de licenciatura.}
\end{table}

\subsubsection*{Énfasis en matemática pura}
\begin{itemize}
    \item \hyperlink{MA0982}{MA-0982 Seminario de estudios dirigidos en matemática pura I}
    \item \hyperlink{MA1081Pura}{MA-1081 Seminario de estudios dirigidos en matemática pura II}
\end{itemize}

\subsubsection*{Énfasis en matemática aplicada}
\begin{itemize}
    \item \hyperlink{MA0984}{MA-0984 Seminario de estudios dirigidos en matemática aplicada I}
    \item \hyperlink{MA1082}{MA-1082 Seminario de estudios dirigidos en matemática aplicada II}
\end{itemize}

\subsection{Cursos optativos}
\subsubsection*{Cursos optativos de núcleo en matemática pura OPT-001}

\begin{itemize}
    \item \hyperlink{MA0788}{MA-0788 Análisis Real II}
    \item \hyperlink{MA0506}{MA-0506 Teoría algebraica de números}
    \item \hyperlink{MA0809}{MA-0809 Teoría analítica de números}
    \item \hyperlink{MA0512}{MA-0512 Teoría de conjuntos}
    \item \hyperlink{MA0525}{MA-0525 Combinatoria}
    \item \hyperlink{MA0528}{MA-0528 Sistemas dinámicos}
    \item \hyperlink{MA0703}{MA-0703 Integración}
    \item \hyperlink{MA0707}{MA-0707 Geometría Diferencial}
    \item \hyperlink{MA0810}{MA-0810 Geometría algebraica I}
    \item \hyperlink{MA0711}{MA-0711 Lógica}
    \item \hyperlink{MA0714}{MA-0714 Optimización}
    \item \hyperlink{MA0755}{MA-0755 Ecuaciones diferenciales parciales}
    \item \hyperlink{MA0804}{MA-0804 Topología algebraica}
    \item \hyperlink{MA0806}{MA-0806 Análisis funcional}
    \item \hyperlink{MA0815}{MA-0815 Análisis armónico}
    \item \hyperlink{MA0820}{MA-0820 Teoría de modelos}
    \item \hyperlink{MA0840}{MA-0840 Probabilidad}
    \item \hyperlink{MA0860}{MA-0860 Teoría de módulos}
    \item \hyperlink{MA0889}{MA-0889 Álgebra conmutativa}
    \item \hyperlink{MA0920}{MA-0920 Ecuaciones en derivadas parciales numéricas}
\end{itemize}

\subsubsection*{Cursos optativos de núcleo en matemática pura OPT-002}
\begin{itemize}
    \item {CA-0304 Herramientas de ciencia de datos II}
    \item {CA-0411 Análisis de Datos II}
    \item {CA-0403 Estadística actuarial II}
    \item {CA-0512 Modelos Lineales}
    \item {CA-0612 Series de Tiempo}
    \item \hyperlink{MA0406}{MA-0406 Introducción a la optimización}
    \item {MA-0477 Ecuaciones Diferenciales Aplicadas}
    \item {MA-0722 Estadística Matemática I}
\end{itemize}

\subsubsection*{Cursos optativos en general OPT-003}


\begin{itemize}
    \item \hyperlink{MA0788}{MA-0788 Análisis Real II}
    \item \hyperlink{MA0506}{MA-0506 Teoría algebraica de números}
    \item \hyperlink{MA0809}{MA-0809 Teoría analítica de números}
    \item \hyperlink{MA0512}{MA-0512 Teoría de conjuntos}
    \item \hyperlink{MA0525}{MA-0525 Combinatoria}
    \item \hyperlink{MA0528}{MA-0528 Sistemas dinámicos}
    \item \hyperlink{MA0703}{MA-0703 Integración}
    \item \hyperlink{MA0707}{MA-0707 Geometría Diferencial}
    \item \hyperlink{MA0810}{MA-0810 Geometría algebraica I}
    \item \hyperlink{MA0711}{MA-0711 Lógica}
    \item \hyperlink{MA0714}{MA-0714 Optimización}
    \item \hyperlink{MA0755}{MA-0755 Ecuaciones diferenciales parciales}
    \item \hyperlink{MA0804}{MA-0804 Topología algebraica}
    \item \hyperlink{MA0806}{MA-0806 Análisis funcional}
    \item \hyperlink{MA0815}{MA-0815 Análisis armónico}
    \item \hyperlink{MA0820}{MA-0820 Teoría de modelos}
    \item \hyperlink{MA0840}{MA-0840 Probabilidad}
    \item \hyperlink{MA0860}{MA-0860 Teoría de módulos}
    \item \hyperlink{MA0889}{MA-0889 Álgebra conmutativa}
    \item \hyperlink{MA0920}{MA-0920 Ecuaciones en derivadas parciales numéricas}
    \item {CA-304 Herramientas de ciencia de datos II}
    \item {CA-0411 Análisis de Datos II}
    \item {CA-0403 Estadística actuarial II}
    \item {CA-0512 Modelos Lineales}
    \item {CA-0612 Series de Tiempo}
    \item \hyperlink{MA0406}{MA-0406 Introducción a la optimización}
    \item {MA-0477 Ecuaciones Diferenciales Aplicadas}
    \item {MA-0722 Estadística Matemática I}    
    \item \hyperlink{MA0809}{MA-0809 Teoría analítica de números}
    \item \hyperlink{MA0790}{MA-0790 Tópicos de análisis}
    \item \hyperlink{MA0791}{MA-0791 Tópicos de Probabilidad}
    \item \hyperlink{MA0792}{MA-0792 Tópicos de Ecuaciones diferenciales}
    \item \hyperlink{MA0793}{MA-0793 Tópicos de Geometría}
    \item \hyperlink{MA0794}{MA-0794 Tópicos de Modelación}
    \item \hyperlink{MA0796}{MA-0796 Tópicos de Álgebra}
    \item \hyperlink{MA0797}{MA-0797 Tópicos de Matemática discreta}
    \item \hyperlink{MA0830}{MA-0830 Tópicos de Teoría de Números}
    \item \hyperlink{MA0831}{MA-0831 Tópicos de Lógica}
    \item \hyperlink{MA0832}{MA-0832 Tópicos en Topología}
    \item \hyperlink{MA0918}{MA-0918 Procesos estocásticos}
    \item \hyperlink{MA0794}{MA-0794 Tópicos de Modelación}
    \item \hyperlink{MA0795}{MA-0795 Tópicos de Computación científica}
    \item \hyperlink{MA0798}{MA-0798 Tópicos de Análisis de datos}
    \item \hyperlink{MA0799}{MA-0799 Tópicos de Aprendizaje automático}
    \item {MA-0759 Tópicos de Matemática Financiera}
    \item \hyperlink{MA0917}{MA-0917 Estadística II}
    \item \hyperlink{MA0918}{MA-0918 Procesos estocásticos}
    \item {CA-917 Estadística matemática II}
    \item \hyperlink{MA0795}{MA-0795 Tópicos de Computación científica}
\end{itemize}


